\begin{exercise}{145}
Show that a map $f\colon K\to L$ is cofinal if and only if $f^{\mathrm{op}}\colon K^{\mathrm{op}}\to L^{\mathrm{op}}$ is
coinitial.

\begin{proof}
Taking opposites is an involution preserving products and pullbacks.
The exponential law therefore gives natural isomorphisms
\[
 \mathrm{Fun}(A,B)^{\mathrm{op}}
 \cong\mathrm{Fun}(A^{\mathrm{op}},B^{\mathrm{op}}).
\]
For a right fibration $p\colon X\to L$, these restrict to
\[
 \begin{aligned}
 \mathrm{Fun}_L(L,X)^{\mathrm{op}}
   &\cong\mathrm{Fun}_{L^{\mathrm{op}}}
       (L^{\mathrm{op}},X^{\mathrm{op}}),\\
 \mathrm{Fun}_f(K,X)^{\mathrm{op}}
   &\cong\mathrm{Fun}_{f^{\mathrm{op}}}
       (K^{\mathrm{op}},X^{\mathrm{op}}),
 \end{aligned}
\]
since the fixed vertices $\mathrm{id}_L$ and $f$ become
$\mathrm{id}_{L^{\mathrm{op}}}$ and $f^{\mathrm{op}}$.
Naturality identifies the opposite of restriction along $f$ with
restriction along $f^{\mathrm{op}}$.

These spaces are Kan complexes by Remark~4.4.3.
Taking opposites preserves homotopy equivalences: apply it to a
homotopy inverse and its homotopies, using
$(\Delta^1)^{\mathrm{op}}\cong\Delta^1$ to reverse the homotopy parameter.
Observation~2.2.11 thus shows that restriction is a Joyal equivalence
exactly when its opposite is.
By Exercise~\ref{ex:51}(3), right fibrations over $L$ correspond
bijectively to left fibrations over $L^{\mathrm{op}}$.
Consequently the two tests in Definition~4.4.2 are equivalent,
which proves both implications.
\end{proof}
\end{exercise}
